\section{Worked kernel: conv2d}
\label{sec:conv2d}

\subsection{The declaration, axis by axis}

\begin{lstlisting}
#[tiled_kernel]
#[tile_loop(trip_count = [C_IN, G, KH, KW],
            axes = [c_in_local = (C_IN / G), kh = KH, kw = KW], generate)]
#[tile_carry(acc = [BLOCK_OW])]
pub fn conv2d_forward<T: Triton, D: Num, const KH: i32, const KW: i32,
    const STRIDE_H: i32, const STRIDE_W: i32, const PAD_H: i32, const PAD_W: i32,
    const G: i32, const BLOCK_OW: i32>(
    #[tile(name = "B", extent = _B)]                                     // dim 0: untiled
    #[tile(extent = C_IN)]                                               // dim 1: untiled
    #[tile(extent = H, window(stride = STRIDE_H, pad = PAD_H, kernel = KH, output = OH))]
    #[tile(block = BLOCK_OW, extent = W,
           window(stride = STRIDE_W, pad = PAD_W, kernel = KW, output = OW))]
    #[tile_loop_tile(index = [tile_b = _B,
        ((tile_c_out / (C_OUT / G)) * (C_IN / G) + c_in_local) = C_IN,
        (tile_oh * STRIDE_H + kh - PAD_H) = H,
        (__tile_range * STRIDE_W + kw - PAD_W) = W], bounds = [2, 3])]
    x: In<Tile<T, D>>,
    #[tile_loop_scalar(index = [tile_c_out = C_OUT, c_in_local = (C_IN / G), kh = KH, kw = KW])]
    w: In<Tile<T, D>>,                                                   // no #[tile]: not in spec
    #[tile(name = "B", extent = _B)]
    #[tile(extent = C_OUT)]
    #[tile(extent = OH)]
    #[tile(block = BLOCK_OW, extent = OW)]
    y: Out<Tile<T, D>>,
    _B: i32, C_IN: i32, C_OUT: i32, H: i32, W: i32, OH: i32, OW: i32,
) { acc = acc + x * w; }
\end{lstlisting}

\begin{itemize}
\item \textbf{\code{x}'s H and W axes} name their own extents (\code{H}, \code{W}) and,
  through \code{output = OH/OW}, the output variable they resolve against. Naming the
  input axis \code{OW} would also resolve, but the spec would then claim the input's
  width \emph{is} \code{OW}, which \code{untiled_dims} and any footprint model would
  believe.
\item \textbf{The input and output axes match through windows.} The prelude check
  (``every input axis matches an output axis'') is satisfied by the windows. \code{x}
  is loop-indexed, so the prelude does not place it anyway.
\item \textbf{\code{#[tile_loop_tile]} restates the coordinates as expressions.} The
  window says \emph{what} is read; the index list says \emph{where}. They have to
  agree, and nothing checks that they do (\S\ref{sec:gotchas}).
\end{itemize}

\subsection{What \texttt{tile\_spec()} and \texttt{grid\_spec()} return}

\begin{lstlisting}
KernelTileSpec {
  inputs: [ TensorTileSpec { param: "x", rank: 4, reduction_axis: None,
              untiled_dims: ["B", "C_IN"],
              axes: [ { dims: [2], block_const: "1",        extent_param: "H",
                        window: Some({ output: "OH", stride: "STRIDE_H", pad: Some("PAD_H"), kernel: "KH" }) },
                      { dims: [3], block_const: "BLOCK_OW", extent_param: "W",
                        window: Some({ output: "OW", stride: "STRIDE_W", pad: Some("PAD_W"), kernel: "KW" }) } ] } ],
  outputs: [ TensorTileSpec { param: "y", rank: 4, reduction_axis: None,
              untiled_dims: ["B", "C_OUT", "OH"],
              axes: [ { dims: [3], block_const: "BLOCK_OW", extent_param: "OW", window: None } ] } ],
  loop_spec: Some(TileLoopSpec { carries: [{ name: "acc", shape_consts: ["BLOCK_OW"] }],
                                 trip_count_factors: ["C_IN", "G", "KH", "KW"] }),
}
GridSpec { axes: [ B(_B), C_OUT, OH, OW(BLOCK_OW) ] }      // all on GridDim::X, in decode order
\end{lstlisting}

\subsection{Resolving it}

\begin{figure}[h]
\centering
\begin{tikzpicture}[x=0.42cm,y=0.42cm,>={Stealth[length=1.8mm]}]
  % input row
  \foreach \i in {0,...,17} { \draw[fill=attr!12,draw=boxline] (\i,0) rectangle ++(1,1); }
  \foreach \i in {-1,18} { \draw[fill=bg,draw=boxline,dashed] (\i,0) rectangle ++(1,1); }
  \node[font=\scriptsize,anchor=east] at (-1.4,0.5) {input \texttt{x} (W)};
  % output row
  \foreach \i in {0,...,15} { \draw[fill=st!18,draw=boxline] (\i+1,3.2) rectangle ++(1,1); }
  \node[font=\scriptsize,anchor=east] at (-1.4,3.7) {output \texttt{y} (OW)};
  % window lines for first and last output
  \draw[color=cm] (1.5,3.2) -- (0.5,1); \draw[color=cm] (1.5,3.2) -- (2.5,1);
  \draw[color=cm] (16.5,3.2) -- (15.5,1); \draw[color=cm] (16.5,3.2) -- (17.5,1);
  \draw[|<->|] (1,4.6) -- (17,4.6) node[midway,above,font=\scriptsize] {\texttt{BLOCK\_OW} $=b=16$};
  \draw[|<->|] (0,-0.5) -- (18,-0.5) node[midway,below,font=\scriptsize] {$(b-1)\cdot s+k = 15\cdot1+3 = 18$};
  \node[font=\scriptsize,text=cm,anchor=west] at (19.3,0.5) {pad: shifts origin,};
  \node[font=\scriptsize,text=cm,anchor=west] at (19.3,-0.2) {never the size};
\end{tikzpicture}
\caption{The receptive field of one conv2d tile along W ($k=3$, $s=1$, $p=1$).}
\label{fig:window}
\end{figure}

\code{test_conv2d_windowed_input_resolves_to_its_receptive_field} takes an output
tile of \code{[2, 16, 5, 8]} with $\text{STRIDE\_W}=1$, $\text{KW}=3$. The walk:

\begin{enumerate}
\item \textbf{Read the output's bindings.} \code{y} has one binding, \code{OW} on
  dim~3, so \code{resolved = {OW: 8}}. The \code{2}, \code{16} and \code{5} sit on
  untiled dims and \emph{are discarded}.
\item \textbf{\code{x} dim 3 (W):} windowed against \code{OW}, so
  $(8-1)\cdot 1+3=\mathbf{10}$. With \code{PAD_W = 0} it is still 10, and with stride 2
  it is 17.
\item \textbf{\code{x} dim 2 (H):} windowed against \code{OH}, but \code{OH} is not in
  \code{resolved}, so it stays \code{None}: all of \code{H}. The receptive-field
  formula would give $\text{KH}$ for the one output row a program computes, but the
  output never says it computes one row.
\item \textbf{\code{x} dims 0, 1 (B, C\_IN):} no binding, so \code{None}.
\end{enumerate}

Result: \code{x} needs \code{[None, None, None, Some(10)]}, which is
$B \times C_{\text{IN}} \times H \times 10$ elements. The weight needs nothing,
because it is not in the spec.

\subsection{How large the over-estimate is}
\label{sec:conv-overestimate}

What one program of the expanded body in \S\ref{sec:prelude} actually reads over its
whole loop: $1 \times C_{\text{IN}}/G \times \text{KH} \times ((b-1)s+\text{KW})$
elements of \code{x}, plus $C_{\text{IN}}/G \times \text{KH}\times\text{KW}$ weights.
What the spec says, for the two shapes used in \S\ref{sec:welder}, with $b=16$:

\begin{center}
\small
\begin{tabular}{lrrr}
\toprule
shape (B, C$_\text{in}$$\to$C$_\text{out}$, H$\times$W) & real \code{x}+\code{w} per program & spec's \code{x} per program & ratio \\
\midrule
test graph (10, 3$\to$8, 32$\times$32) & 756 B & 69{,}120 B & 91$\times$ \\
ResNet stage (32, 64$\to$64, 56$\times$56) & 16{,}128 B & 8{,}257{,}536 B & 512$\times$ \\
\bottomrule
\end{tabular}
\end{center}

Against a 48~KB shared-memory budget, the real tile fits easily and the declared one
never does. A Welder footprint check run on today's spec would reject every conv2d
tile at the shared level. Three independent gaps cause this:

\begin{enumerate}
\item \textbf{Grid-decoded output axes have no binding}, so the output tile's value
  there cannot seed anything. A program covers exactly one index of \code{B},
  \code{C_OUT} and \code{OH}. That is a block of 1, as real as \code{BLOCK_OW}.
  \emph{Proposal:} emit an unblocked, gridded output axis as a binding with
  \code{block_const: "1"}, as is already done for an unblocked \emph{windowed input}
  axis. Then \code{B} resolves to 1, and \code{H} resolves through its window to
  $(1-1)\cdot\text{STRIDE\_H}+\text{KH} = \text{KH}$. The test comment ``H unresolved,
  so full extent --- not an over-read'' would no longer hold.
\item \textbf{\code{C_IN} vs \code{C_OUT}:} a program reads one \emph{group's}
  \code{C_IN / G} channels, chosen by \code{tile_c_out}. No name relates the two, so
  even with fix~1 \code{C_IN} resolves to full extent. That is correct for $G=1$ and
  $G\times$ too large for grouped or depthwise conv. \code{divide_by} is the field
  meant for this, but the macro cannot emit it per instance (\issue{1nr.15}).
\item \textbf{The weight is invisible.} Its traffic ($C_{\text{IN}}/G \cdot K^2$ per
  program, re-read by every program) never reaches a cost model. Naming its axes
  needs an \code{extent} that can be an expression (\code{C_IN / G}) or a dedicated
  per-group param.
\end{enumerate}

\begin{note}[What propagation can and cannot vary]
The only searchable dimension of \code{conv2d_forward} is \code{BLOCK_OW}. A tile
with several output rows, several output channels, or a chunk of the
\code{C_IN}$\cdot$\code{KH}$\cdot$\code{KW} reduction is not a different
\emph{configuration} of this kernel. It is a different kernel. \S\ref{sec:welder}
shows that this, rather than the fusion question, decides most of the traffic.
\end{note}

\subsection{The tiled variant: two blocked axes}

\code{conv2d_bn_silu_tiled_forward} (route~2, hand-written body) computes a
\code{[BLOCK_N, BLOCK_OW]} block of (output channels $\times$ columns) using an
outer-product update:

\begin{lstlisting}
#[tile_loop(trip_count = [C_IN, KH, KW])]
#[tile_carry(acc = [BLOCK_N, BLOCK_OW])]
#[tile_grid(order = [B, OH, C_OUT, OW])]          // body decodes (b, oh, n_tile, ow_tile)
    ...
    #[tile(name = "B", extent = B)]
    #[tile(block = BLOCK_N, extent = C_OUT)]
    #[tile(extent = OH)]
    #[tile(block = BLOCK_OW, extent = OW)]
    y_ptr: Out<T::Pointer<f32>>,
\end{lstlisting}

Its \code{y} carries two bindings, \code{C_OUT} by \code{BLOCK_N} and \code{OW} by
\code{BLOCK_OW}. An output tile \code{[_, 32, _, 16]} therefore seeds both. \code{x}
still declares no \code{C_IN} binding, because every program reads all input
channels, which is correct here since $G=1$ is enforced at dispatch. The extra
binding is what makes the channel-tiled shape a configuration the search can reach.
The grid order is stated because the decode is not the output's dim order. A rider
fused under this grid would otherwise be handed the wrong indices.

\code{conv2d_bn_silu_gemm_forward} (implicit GEMM for $1\times1$ convs) is
deliberately \emph{undeclared}. Its \code{M = OH * OW} is a flattened spatial axis,
which \code{span} could now express. But its output is written as
\code{[B * C_OUT, y_row_stride]}, with a TMA-aligned allocation stride wider than
\code{M}, and its \code{GROUP_M} decode would need \code{swizzled}.
